Tool-augmented data agents can run a tool successfully and still reach a wrong conclusion when its output is silently corrupted. We introduce ToxicBench, a benchmark that inserts numerical, label, schema, and retrieval errors into otherwise successful tool responses. Each task pairs clean and poisoned runs with the same query and source data. The benchmark includes 120 single-table instances, a 13-task multi-table extension, and a cross-model suite with 34 numerical and 24 semantic/schema instances. Under automatic scoring, clean task success of 0.73--0.90 falls to 0.49--0.53 in the expanded GPT evaluation. Trajectory metrics separate poisoned-answer adoption, evidence checking, and recovery. Two annotators evaluate 240 held-out trajectories to test both the scoring decisions and method comparisons. Matched-budget protocols compare an ordinary second pass with explicit verification, while trace analysis records whether later observations remain clean. Under repeated poisoning, both annotators identify cases that check further evidence and still adopt the poisoned conclusion. ToxicBench provides controlled task pairs and an auditable evaluation protocol for studying when agents trust tool evidence and when additional checks help.
